\section{文档介绍}

\subsection{文档目的}

本文档为《学术图像造假检测系统》的需求规格说明书，旨在明确系统开发的目标、功能与非功能性需求，作为开发团队设计、实现和测试的依据，并为后续项目验收、用户反馈及迭代优化提供基准。

本文档面向软件开发人员和使用本产品的全部用户，希冀进一步规范软件开发的细节工作，并且全方位展示软件项目的具体实施情况，包括但不限于项目概述、功能性和非功能性需求的阐述、外围系统和接口要求等方面的进展。

\subsection{文档范围}

该份需求规格说明书主要包括：

\begin{itemize}
    \item 本文档应提前说明并列示出其编写目的及通篇架构，同时标明其参考文档和专业性术语。
    \item 本文档使用系统和软件项目的概述情况，包括其基本标识如名称、缩略名和标识号等，同样包括项目的背景、目标、整体结构性框架，对用户的群像描述及设计层面的设想与阻碍。
    \item 本文档所适用的软件用途（即软件开发人员所设想的全部需求功能）。它还应描述系统与软件的一般特性；概述系统开发、运行和维护的相关设想；标识项目的需方、用户、开发方等；如有需求，列出其他有关文档。
\end{itemize}

\subsection{读者对象}

本文档主要面向以下读者：

\begin{itemize}
    \item 开发团队：负责系统设计、开发和测试的人员。
    \item 项目管理人员：负责项目进度、质量控制和资源调配的人员。
    \item 客户与最终用户：使用该学术图像造假检测系统的学术研究人员、教育机构、出版机构等。
\end{itemize}

\subsection{参考文档}

\begin{center}
\renewcommand{\arraystretch}{1.5}
\begin{tabular}{|p{1.5cm}|p{6cm}|p{3cm}|p{2.5cm}|}
\hline
\centering 编号 & \centering 文档名称 & \centering 单位 & \centering 日期 \tabularnewline
\hline
\centering 1 & A3P-FR-3\_需求规格书\_0327 & \centering 北京航空航天大学 & \centering 2024-03-27 \tabularnewline
\hline
\end{tabular}
\end{center}

\subsection{术语与缩写解释}

\begin{center}
\renewcommand{\arraystretch}{1.5}
\begin{longtable}{|p{3.5cm}|p{9cm}|}
\hline
\centering 缩写、术语 & \centering 解 释 \tabularnewline
\hline
Web端 & 开发人员对电脑端的网页版（通常从视觉上呈现给用户）进行的开发管理。 \tabularnewline
\hline
APP & Application的缩写，主要指安装在智能手机上的软件，完善原始系统的不足与个性化。 \tabularnewline
\hline
BIB & BibTeX的缩写，是一种文件格式，也可指代制作BibTex文件的工具。这种文件用于描述和处理引用列表，通常情况下与 LaTeX 文档结合使用。 \tabularnewline
\hline
AI 图像检测算法 & 基于深度学习模型（如CNN、GAN识别算法）对图像进行篡改痕迹分析（如复制-粘贴、重复使用、PS痕迹等）。 \tabularnewline
\hline
Chrome/360浏览器 /Windows10 & 用来检索、展示以及传递Web信息资源的不同应用程序。 \tabularnewline
\hline
Django & Python（一种编程语言）驱动的开源模型—视图—控制图风格的Web应用程序框架，可以在短时间内创建高质量、易维护、数据库驱动的应用程序。 \tabularnewline
\hline
UI & 用户界面（User Interface）是指对软件的人机交互、操作逻辑、界面美观的整体设计。 \tabularnewline
\hline
时间戳 & 使用数字签名技术产生的数据，签名的对象包括了原始文件信息、签名参数、签名时间等信息。 \tabularnewline
\hline
\end{longtable}
\end{center}